\begin{theorem}[Jensen’s Inequality]
Let \(f:\mathbb{R}\to\mathbb{R}\) be convex, i.e.  
\[
\forall x,y\in\mathbb{R},\;\forall t\in[0,1]\;:\qquad 
f\bigl(tx+(1-t)y\bigr)\le t\,f(x)+(1-t)\,f(y). \tag{C}
\]
Let \(x_{1},\dots ,x_{n}\in\mathbb{R}\) and let weights \(\lambda_{1},\dots ,\lambda_{n}\) satisfy  
\[
\lambda_{i}\ge 0\;(i=1,\dots ,n),\qquad \sum_{i=1}^{n}\lambda_{i}=1 .
\]
Then  
\[
f\!\left(\sum_{i=1}^{n}\lambda_{i}x_{i}\right)\le 
\sum_{i=1}^{n}\lambda_{i}f(x_{i}). \tag{J}
\]
\end{theorem}

\begin{proof}
\textbf{1. Degenerate cases.}  

If \(\lambda_{k}=1\) and \(\lambda_{i}=0\) for \(i\neq k\), then  
\(\sum_{i=1}^{n}\lambda_{i}x_{i}=x_{k}\) and 
\(\sum_{i=1}^{n}\lambda_{i}f(x_{i})=f(x_{k})\); (J) reduces to the identity  
\(f(x_{k})\le f(x_{k})\).  

If some weight equals \(0\), the corresponding index can be omitted from both
sides of (J). Hence we may assume that at least two weights are positive.

\textbf{2. Binary case \((n=2)\).}  

Let \(\lambda_{1},\lambda_{2}\ge0\) with \(\lambda_{1}+\lambda_{2}=1\).  
Applying (C) with \(t=\lambda_{1}\), \(x=x_{1}\), \(y=x_{2}\) yields  
\[
f\bigl(\lambda_{1}x_{1}+\lambda_{2}x_{2}\bigr)
\le \lambda_{1}f(x_{1})+\lambda_{2}f(x_{2}),
\]
which is (J) for \(n=2\).

\textbf{3. Inductive step.}  

Assume (J) holds for any collection of \(n-1\) points (\(n\ge3\)).
Let \(x_{1},\dots ,x_{n}\) and non‑negative weights \(\lambda_{1},\dots ,\lambda_{n}\) satisfy \(\sum_{i=1}^{n}\lambda_{i}=1\).

Set  
\[
L:=\sum_{i=1}^{n-1}\lambda_{i},\qquad 1-L=\lambda_{n}.
\]
Since the \(\lambda_{i}\) are non‑negative, \(0\le L\le1\).

If \(L=0\) (resp.\ \(L=1\)), we are back to the trivial single‑weight case or to the
induction hypothesis applied to the first \(n-1\) points.  
Thus we may assume \(0<L<1\).

Define normalized weights for the first \(n-1\) points:
\[
\mu_{i}:=\frac{\lambda_{i}}{L}\qquad (i=1,\dots ,n-1),
\]
so that \(\mu_{i}\ge0\) and \(\sum_{i=1}^{n-1}\mu_{i}=1\).  
Let  
\[
y:=\sum_{i=1}^{n-1}\mu_{i}x_{i}.
\]

\textit{Induction hypothesis} applied to \(\{x_{i}\}_{i=1}^{n-1}\) gives  
\begin{equation}
f(y)\le\sum_{i=1}^{n-1}\mu_{i}f(x_{i}). \tag{IH}
\end{equation}

The full convex combination can be rewritten as  
\begin{equation}
\sum_{i=1}^{n}\lambda_{i}x_{i}=L\,y+(1-L)\,x_{n}. \tag{*}
\end{equation}

Applying convexity (C) to the pair \((y,x_{n})\) with parameter \(t=L\) yields  
\begin{equation}
f\bigl(Ly+(1-L)x_{n}\bigr)\le L\,f(y)+(1-L)f(x_{n}). \tag{C2}
\end{equation}

Using (IH) and the definition of \(\mu_{i}\),
\[
L\,f(y)=L\sum_{i=1}^{n-1}\mu_{i}f(x_{i})
       =\sum_{i=1}^{n-1}\lambda_{i}f(x_{i}). \tag{*1}
\]

Now
\[
\begin{aligned}
f\!\left(\sum_{i=1}^{n}\lambda_{i}x_{i}\right)
&\overset{(*)}{=}f\bigl(Ly+(1-L)x_{n}\bigr) \\
&\le L\,f(y)+(1-L)f(x_{n}) \\
&\overset{(*1)}{=}\sum_{i=1}^{n-1}\lambda_{i}f(x_{i})+\lambda_{n}f(x_{n}) \\
&=\sum_{i=1}^{n}\lambda_{i}f(x_{i}),
\end{aligned}
\]
which is precisely (J) for \(n\) points.  
Thus, by induction, (J) holds for every integer \(n\ge1\).

\textbf{4. Equality condition.}  

In the binary inequality (C) equality holds iff either \(t=0\) or \(t=1\) (one
weight is zero) or the restriction of \(f\) to the segment \([x,y]\) is affine.
In the inductive proof we used (C) once, with \(t=L\) and points \(y\) and
\(x_{n}\). Hence equality in (J) occurs precisely when  

\begin{enumerate}
\item \(L=0\) or \(L=1\) (all weight concentrated on a single point), or
\item \(f\) is affine on the segment \([y,x_{n}]\) and equality already held in
the induction hypothesis for the first \(n-1\) points.
\end{enumerate}

Iterating this observation shows that equality holds iff either all but one
weight are zero, or \(f\) is affine on the convex hull of the points with
positive weight. In particular, if all such points coincide or if \(f\) is
globally affine (\(f(t)=at+b\)), equality holds for every choice of
\(\{x_{i}\},\{\lambda_{i}\}\).

\end{proof}